\section*{Sets and Indices}

\begin{itemize}
    \item $I = \{1,\dots,n\}$: item types (customer-required pipe lengths), indexed by $i$.
    \item $P = \{1,\dots,|P|\}$: set of all feasible cutting patterns generated by enumeration, indexed by $p$.
    \item $K = \{1,\dots,\bar K\}$: frequency-rank positions for pattern-usage surcharge, indexed by $k$.
\end{itemize}

For this instance:
\[
I=\{1,2,3,4\},\qquad \bar K = 4.
\]

\section*{Parameters}

\begin{itemize}
    \item $L^{\text{raw}}$ (\texttt{raw\_pipe\_length}): length of each raw pipe. Here $L^{\text{raw}}=1850$.
    \item $\bar K$ (\texttt{max\_cutting\_patterns}): maximum number of distinct cutting patterns allowed. Here $\bar K=4$.
    \item $C^{\max}$ (\texttt{max\_cuts\_per\_pipe}): maximum number of pieces cut from one raw pipe. Here $C^{\max}=5$.
    \item $W^{\max}$ (\texttt{max\_leftover\_length}): maximum allowed leftover length in any used pattern. Here $W^{\max}=100$.
    \item $L^{\min}=L^{\text{raw}}-W^{\max}$ (\texttt{min\_length}): minimum utilized length for any feasible pattern. Here $L^{\min}=1750$.
    \item $\ell_i$: required finished length of item type $i$.
    \item $d_i$: demand quantity of item type $i$.
    \item $r_k$: surcharge rate associated with using at least $k$ distinct patterns, in rank order. From the data,
    \[
    r_1=0.1,\quad r_2=0.2,\quad r_3=0.3,
    \]
    and no additional rate is provided for $k=4$.
    \item $a_{pi}$: number of pieces of item type $i$ produced by pattern $p$.
    \item $M=\sum_{i\in I} d_i$: big-$M$ constant used in the linking constraints. For this instance,
    \[
    M=15+28+21+30=94.
    \]
\end{itemize}

Demand data:
\[
(\ell_i,d_i)\in\{(290,15),(315,28),(350,21),(455,30)\}.
\]

The pattern set $P$ is the set of all integer vectors $a_p=(a_{pi})_{i\in I}$ such that
\[
\sum_{i\in I} a_{pi}\ell_i \in [L^{\min},L^{\text{raw}}],
\qquad
1 \le \sum_{i\in I} a_{pi} \le C^{\max},
\qquad
a_{pi}\in \mathbb{Z}_{\ge 0}.
\]
Equivalently, each enumerated pattern has total used length between $1750$ and $1850$ and contains at most $5$ cut pieces.

\section*{Decision Variables}

\begin{itemize}
    \item $y_p$ (code: \texttt{y[p]}): integer number of raw pipes cut according to pattern $p$.
    \item $z_p$ (code: \texttt{z[p]}): binary indicator equal to $1$ if pattern $p$ is used, and $0$ otherwise.
    \item $u_k$ (code: \texttt{u[k]}): binary indicator equal to $1$ if at least $k$ distinct patterns are used, and $0$ otherwise.
\end{itemize}

Domains:
\[
y_p\in \mathbb{Z}_{\ge 0}\quad(\forall p\in P),\qquad
z_p\in\{0,1\}\quad(\forall p\in P),\qquad
u_k\in\{0,1\}\quad(\forall k\in K).
\]

\section*{Objective}

The implemented model minimizes the total number of raw pipes used plus the rank-based surcharge on the number of distinct patterns:
\[
\min \sum_{p\in P} y_p \;+\; \sum_{k\in K:\,k\le |r|} r_k\,u_k,
\]
where $|r|$ is the number of surcharge rates present in the parameter file. For this instance,
\[
\min \sum_{p\in P} y_p + 0.1\,u_1 + 0.2\,u_2 + 0.3\,u_3.
\]

\section*{Constraints}

\paragraph{Demand satisfaction.}
For each required length, total production across all chosen patterns must meet or exceed demand:
\[
\sum_{p\in P} a_{pi}\,y_p \ge d_i
\qquad \forall i\in I.
\]

\paragraph{Pattern-use linking.}
A pattern can produce pipes only if it is declared used:
\[
y_p \le M z_p
\qquad \forall p\in P.
\]

\paragraph{Counting used patterns via rank variables.}
The number of used patterns equals the number of activated rank positions:
\[
\sum_{p\in P} z_p = \sum_{k\in K} u_k.
\]

\paragraph{Monotonicity of rank variables.}
If the $(k+1)$-st rank is active, then the $k$-th rank must also be active:
\[
u_k \ge u_{k+1}
\qquad \forall k=1,\dots,\bar K-1.
\]

\paragraph{Maximum number of distinct patterns.}
\[
\sum_{p\in P} z_p \le \bar K.
\]

\section*{Notes on Implementation}

\begin{itemize}
    \item The code first enumerates the complete finite set of feasible cutting patterns $P$ recursively. No pattern-generation/column-generation scheme is used.
    \item A feasible pattern is any nonzero integer composition of the item lengths such that:
    \[
    1750 \le \sum_i a_{pi}\ell_i \le 1850,
    \qquad
    \sum_i a_{pi}\le 5.
    \]
    This exactly reflects the code logic using \texttt{min\_length = raw\_pipe\_length - max\_leftover}.
    \item The surcharge variables $u_k$ do not identify which pattern is most frequent, second most frequent, etc. Instead, they only count how many distinct patterns are used, and the objective adds the first few rates accordingly. This is the exact implemented logic and should not be reinterpreted as a true frequency-ordering model.
    \item Because only three surcharge-rate parameters are present in the data, the fourth rank variable $u_4$ has no direct cost coefficient, although it still participates in the counting and monotonicity constraints.
    \item Demands are allowed to be over-satisfied because the model uses $\ge$ demand constraints and has no upper bounds on produced quantities other than those induced by minimizing the number of raw pipes.
    \item The optimization problem solved in the code is therefore a finite mixed-integer linear program over the enumerated pattern set, solved by Gurobi through PuLP compatibility.
\end{itemize}
